% FILENAME: paper/section-Matter.tex

\section{Matter: Topological Knots and the SU(3) Problem}

\subsection{The Necessity of Non-Abelian Fields}
% TODO for Author: Explain why a U(1) Abelian field (like pure electromagnetism) cannot create space on its own. The volume determinant natively collapses. The universe absolutely requires non-Abelian fields (color) to expand into a macroscopic geometry.

\begin{equation}
    [F_{\mu\nu}, F_{\rho\sigma}] = 0 \implies \det(g_{\mu\nu}) = 0
\end{equation}
\begin{minted}{lean}
-- CGD.Particles.Color
theorem kinematicSingleColorDegeneracy :
  ∀ (F : Fin 4 → Fin 4 → SL2C),
    isSingleColor F →
    (urbantkeMetric F).det = 0
\end{minted}

\subsection{Topological Stability}
% TODO for Author: Explain that particles are not point masses added to the universe; they are stable topological knots (instantons) in the gauge field. Introduce the Cartan-Maurer degree / Belavin bounds to explain why these knots cannot mathematically decay or unwrap into the trivial vacuum.

\begin{equation}
    \int_{S^3} \text{Cartan-Maurer Form} \neq 0 \implies A \text{ cannot decay to } 0
\end{equation}
\begin{minted}{lean}
-- CGD.Particles.TopologicalStability
theorem kinematicTopologicalStability 
  (windingNumber : (S3 → SU2Group) → ℤ)
  (cartanMaurerIntegral : (S3 → SU2Group) → ℝ)
  ... [Litlib.Belavin1975 and Nakahara Cartan-Maurer boundaries elided] ...
  (integral_zero : cartanMaurerIntegral 1 = 0) :
  ¬ isHomotopicConnection bpstInstanton 0
\end{minted}

\subsection{Geometric Confinement}
% TODO for Author: Address the "SU(3) Problem". Why don't we need a fundamental SU(3) gauge group for the Strong Force? Explain that multi-color components natively crush into 1D topological flux tubes (strings), perfectly yielding a 1D confining Hamiltonian.

\begin{equation}
    \mathcal{H}_{3D} \rightarrow \frac{1}{2} E_z^2 + \frac{1}{2} B^2 \quad \text{(1D Confinement Reduction)}
\end{equation}
\begin{minted}{lean}
-- CGD.Particles.Confinement
theorem kinematicStringConfinement :
  ∀ (E B : Matrix (Fin 3) (Fin 3) ℂ) (E_z : ℂ),
    isCrushedString E E_z →
    ∃ (v : Fin 3 → ℂ), (∑ i : Fin 3, v i * v i = 1) ∧
      densitizedHamiltonian E B = (1 / 2 : ℂ) * E_z^2 + (1 / 2 : ℂ) * ∑ a : Fin 3, ∑ b : Fin 3, B a b * B a b
\end{minted}
